\documentclass{article}
\usepackage[utf8]{inputenc}
\usepackage{a4wide}

\begin{document}

\section*{miniSolitaire}

พิจารณาเกมโซลิแทร์ที่มีไพ่ 3 กองจำนวน\texttt{x,\ y,\ z}ใบตามลำดับ การเล่นแต่ละครั้ง

ดร.ต๋อยสามารถเลือก 2 กองจากกองที่ยังเหลือไพ่อยู่ เพื่อที่จะดึงออกมากองละ 1 ใบ

และเพิ่มคะแนน 1 แต้ม\textbf{เกมจะหยุดเมื่อเหลือไพ่ไม่เกิน 1

กองที่ไม่ว่าง}(นั่นคือจะเล่นต่อไม่ได้)\\



\hfill\break



\textbf{ตัวอย่างที่ 1}



\begin{quote}

\textbf{Input:}2 4 6



\textbf{Output:}6



เริ่มที่ 2 4 6 สามารถเล่นได้วิธีหนึ่งดังนี้\\



ครั้งที่ 1 ดึง 1 3 เหลือ 1 4 5



ครั้งที่ 2 ดึง 1 3 เหลือ 0 4 4



ครั้งที่ 3 ดึง 2 3 เหลือ 0 3 3



ครั้งที่ 4 ดึง 2 3 เหลือ 0 2 2



ครั้งที่ 5 ดึง 2 3 เหลือ 0 1 1



ครั้งที่ 6 ดึง 2 3 เหลือ 0 0 0

\end{quote}



{\textbf{ตัวอย่างที่ 2}}\\



\begin{quote}

\textbf{Input:}4 4 6



\textbf{Output:}7



เริ่มที่ 4 4 6 สามารถเล่นได้วิธีหนึ่งดังนี้\\



ครั้งที่ 1 ดึง 1 2 เหลือ 3 3 6



ครั้งที่ 2 ดึง 1 3 เหลือ 2 3 5



ครั้งที่ 3 ดึง 1 3 เหลือ 1 3 4



ครั้งที่ 4 ดึง 1 3 เหลือ 0 3 3



ครั้งที่ 5 ดึง 2 3 เหลือ 0 2 2



ครั้งที่ 6 ดึง 2 3 เหลือ 0 1 1



ครั้งที่ 7 ดึง 2 3 เหลือ 0 0 0

\end{quote}



{\textbf{\hfill\break

}}



\subsection*{Input}
บรรทัดเดียว เลขจำนวนเต็ม 3 ค่า\texttt{xyz}ตามลำดับ เว้นด้วยช่องว่าง 1 ช่อง\\



\subsection*{Output}
{\textbf{คะแนนสูงสุด}}ที่ดร.ต๋อยสามารถทำได้\\



\end{document}
